\documentclass[11pt,letterpaper]{article}
\usepackage[T1]{fontenc}
\usepackage{lmodern,amsmath,amssymb,amsthm,microtype}
\usepackage[margin=1in]{geometry}
\usepackage[hidelinks]{hyperref}
\newtheorem{theorem}{Theorem}
\title{Three-point extensions of \(G_{449}\) are five-colorable}
\author{Aafi Khan Malek\\ \small Independent researcher, Ahmedabad, India\\
\small\href{mailto:aafimalek2032@gmail.com}{aafimalek2032@gmail.com}}
\date{28 November 2026}
\begin{document}
\maketitle
Let \(P\) be the 449 points given by Heule and Przybocki in their
\href{https://arxiv.org/abs/2611.28518}{\(G_{449}\) construction}.
We use exact finite arithmetic and checked color witnesses to show
that no extension by at most three new points gives a six-chromatic
graph. This does not bound extensions with more points or change
the unrestricted plane bounds.

\begin{theorem}
For any set \(X\) of at most three points in the Euclidean plane,
the induced unit-distance graph on \(P\cup X\) is five-colorable.
\end{theorem}

\section{Points outside the grid}

Let \(L=\tfrac17\mathbb Z[\zeta]\), with \(\zeta=e^{2\pi i/42}\).
The preceding
\href{https://github.com/Aafimalek/novel-discoveries/tree/74e7bf325ddbe52d125453683083c073f451c3fc/hadwiger-nelson/verification/g449/prime-seven-catalogue}
{direction-catalogue proof}
gives all 546 signed unit displacements in \(L\). Applying them to
each base vertex gives 229,320 outside neighbors in \(L\).
Including original vertices as potential grid centers, there are
97,650 distinct triples of base vertices with a common grid center.

For a base triple write its points as \(A/7,B/7,C/7\), and put
\(u=B-A\), \(v=C-A\), \(U=u\bar u\), \(V=v\bar v\),
\(W=(u-v)\overline{(u-v)}\), and \(D=u\bar v-\bar u v\).
The unit circumradius identity is \(UVW+49D^2=0\). A checker
tests this necessary equality at reciprocal primitive
42nd roots modulo \(43,127,211,337,379,463,547\), on all
14,985,824 unordered triples. Every survivor has a checked
grid center. Three distinct points on a circle are noncollinear
and have a unique circumcenter. Thus an off-grid point has at
most two base neighbors.

\section{Adding two points}

The preceding catalogue proof supplies a valid coloring when a
single arbitrary point is added. When a second point has at most
three base neighbors, color it last; it has at most four colored
neighbors. This handles every off-grid second point. The remaining
grid points with at least four base neighbors number 1,043.

For each of their 543,403 pairs, test a bank of 35 proper base
colorings: the supplied coloring and the 34 previously checked
one-point witnesses restricted to the base. For a coloring,
let each added point's list consist of colors absent from its
base neighbors. The pair can be colored if both lists are
nonempty and, when the points are adjacent, their union contains
at least two colors. The original coloring misses 34,911 pairs;
all pairs pass with the bank.

\section{Adding three points}

If one added point has at most two base neighbors, color the
other two first. The last point has at most four colored
neighbors. Consequently only triples from the 2,429 grid points
with at least three base neighbors need further consideration.
List-color the subgraph on the three new points as follows.

\begin{itemize}
\item If there are no new edges, a vertex with at most four base
neighbors can be colored after the other two. Otherwise test
nonempty lists for all three. Failure in the original coloring
must involve one of its 34 one-point exceptions. Grouping the
other candidates by the set of base colorings where their lists
are nonempty gives an exact check.
\item With one edge, first test whether the edge can be colored.
An isolated point with at most four base neighbors causes no
additional condition. For the higher-degree isolated candidates
check that their lists are nonempty in a coloring that handles the edge.
\item In a path, both leaf lists must be nonempty, and a center
color is forbidden precisely when one of the leaves has that
color as its only choice. In a triangle, Hall's condition applies
to the three lists: every list is nonempty, every pairwise union
has size at least two, and the total union has size at least three.
\end{itemize}

The candidate graph has 7,777 edges, 50,995 induced paths and
721 triangles. The bank leaves 60 exceptional triples: 21 with no
edges, 38 with one edge, and one path. Their 54 distinct points
together with \(P\) form a graph on 503 distinct points. A single
five-color vector covers that entire induced graph and hence all
60 cases. A second checker verifies its 2,545 unit edges by
recomputing all 126,253 distances via exact Laurent-polynomial
reduction; it also validates the list rules on all possible masks.
This completes the finite verification.

\paragraph{Scope and reproducibility.}
The companion scripts use pinned input hashes, exact integer
coordinates and an explicit color vector. The replay needs
no SAT solve or floating-point decision. It has not been
formalized in Lean. This is a result about extensions of one
specific graph, not a proof about arbitrary plane configurations.

\paragraph{Contributions.}
Aafi Khan Malek selected and directed the research. ChatGPT
assisted with proof development, software and writing. Codex was
not used. This note has not been peer reviewed.
\end{document}
